% !TEX program = xelatex
\documentclass[11pt,a4paper]{article}

\usepackage[a4paper,margin=18mm,headheight=15pt]{geometry}
\usepackage{fontspec}
\usepackage{polyglossia}
\setdefaultlanguage{russian}
\setotherlanguage{english}
\setmainfont{Arial}
\setsansfont{Arial}
% Arial contains Cyrillic glyphs and keeps inline code readable in Russian text.
\setmonofont{Arial}

\usepackage{microtype}
\usepackage{enumitem}
\usepackage{array}
\usepackage{booktabs}
\usepackage{longtable}
\usepackage{tabularx}
\usepackage{multicol}
\usepackage{xcolor}
\usepackage{hyperref}
\usepackage{fancyhdr}
\usepackage{titlesec}
\usepackage{tcolorbox}
\usepackage{amssymb}
\usepackage{needspace}
\usepackage{lastpage}

\definecolor{MaxPurple}{HTML}{641CF3}
\definecolor{MaxBlue}{HTML}{0A84FF}
\definecolor{MaxCyan}{HTML}{61D7E8}
\definecolor{MaxPink}{HTML}{FF3B80}
\definecolor{Ink}{HTML}{16181D}
\definecolor{Muted}{HTML}{5E6673}
\definecolor{Panel}{HTML}{F3F6FA}
\definecolor{Line}{HTML}{DDE3EC}
\definecolor{PZero}{HTML}{D92D20}
\definecolor{POne}{HTML}{B54708}
\definecolor{PTwo}{HTML}{1570EF}

\hypersetup{
  colorlinks=true,
  linkcolor=MaxPurple,
  urlcolor=MaxBlue,
  pdftitle={План проекта QR-посещаемости в MAX},
  pdfauthor={Команда проекта}
}

\pagestyle{fancy}
\fancyhf{}
\fancyhead[L]{\small\color{Muted}QR-посещаемость в MAX}
\fancyhead[R]{\small\color{Muted}Рабочий план команды}
\fancyfoot[L]{\small\color{Muted}Хакатон: образовательные решения}
\fancyfoot[R]{\small\color{Muted}\thepage\ / \pageref{LastPage}}
\renewcommand{\headrulewidth}{0.4pt}
\renewcommand{\footrulewidth}{0pt}

\titleformat{\section}
  {\Large\bfseries\color{Ink}}
  {\colorbox{MaxPurple}{\textcolor{white}{\strut\thesection}}}
  {0.65em}{}
\titleformat{\subsection}
  {\large\bfseries\color{Ink}}
  {\thesubsection}{0.55em}{}
\titleformat{\subsubsection}
  {\normalsize\bfseries\color{MaxPurple}}
  {\thesubsubsection}{0.5em}{}
\titlespacing*{\section}{0pt}{2.1ex plus .4ex}{1.1ex}
\titlespacing*{\subsection}{0pt}{1.6ex plus .3ex}{0.7ex}

\setlist[itemize]{leftmargin=1.55em,itemsep=0.18em,topsep=0.35em,parsep=0pt}
\setlist[enumerate]{leftmargin=1.7em,itemsep=0.25em,topsep=0.35em}
\newlist{tasklist}{itemize}{2}
\setlist[tasklist]{label=$\square$,leftmargin=1.6em,itemsep=0.16em,topsep=0.3em,parsep=0pt}
\setlist[tasklist,2]{label=$\circ$,leftmargin=1.5em}

\newcommand{\PZero}{\textcolor{PZero}{\textbf{P0}}}
\newcommand{\POne}{\textcolor{POne}{\textbf{P1}}}
\newcommand{\PTwo}{\textcolor{PTwo}{\textbf{P2}}}
\newcommand{\done}[1]{\textbf{Критерий готовности:} #1}
\newcommand{\owner}[1]{\textcolor{Muted}{\textbf{Ответственный:} #1}}
\newcommand{\code}[1]{\texttt{#1}}

\newtcolorbox{summarybox}[1][]{
  colback=Panel,
  colframe=MaxPurple,
  boxrule=0.9pt,
  arc=2mm,
  left=3mm,right=3mm,top=2.5mm,bottom=2.5mm,
  #1
}
\newtcolorbox{warningbox}[1][]{
  colback=white,
  colframe=MaxPink,
  boxrule=0.9pt,
  arc=2mm,
  left=3mm,right=3mm,top=2.5mm,bottom=2.5mm,
  #1
}

\begin{document}

\begin{titlepage}
  \pagecolor{MaxPurple}
  \color{white}
  \vspace*{18mm}
  {\fontsize{15}{18}\selectfont ХАКАТОН \textbullet{} ОБРАЗОВАТЕЛЬНЫЕ РЕШЕНИЯ}\par
  \vspace{22mm}
  {\fontsize{34}{40}\selectfont\bfseries План проекта\\QR-посещаемости\\в MAX}\par
  \vspace{8mm}
  {\Large Полный командный backlog: от исследования проблемы до сдачи и защиты}\par
  \vfill
  \begin{tcolorbox}[colback=white,colframe=white,arc=3mm,boxrule=0pt]
    \color{Ink}
    \textbf{Назначение документа}\par
    Единый рабочий план команды с приоритетами, ответственными, критериями готовности, требованиями к MVP, интеграции с MAX, безопасности, Docker, документации и презентации.
  \end{tcolorbox}
  \vspace{10mm}
  {\large Версия: \today}\par
\end{titlepage}
\nopagecolor
\color{Ink}

\tableofcontents
\newpage

\section*{Как пользоваться документом}
\addcontentsline{toc}{section}{Как пользоваться документом}

\begin{summarybox}
\begin{tabularx}{\textwidth}{>{\bfseries}p{18mm}X}
\PZero & Обязательно для допуска, рабочего MVP и воспроизводимой проверки.\\
\POne & Важно для высокого балла и убедительного пилота.\\
\PTwo & Делать только после стабильного основного сценария.\\
$\square$ & Невыполненная задача. После завершения чекбокс можно отметить вручную.\\
\end{tabularx}
\end{summarybox}

\begin{warningbox}
\textbf{Главное правило проекта.} Пока полностью не работает путь
\emph{«преподаватель запустил сессию $\rightarrow$ студент отметил присутствие $\rightarrow$ преподаватель получил итоговый список»}, команда не переходит к сложной аналитике, внешним интеграциям и платформенному бонусу.
\end{warningbox}

\section{Управление проектом}
\owner{лидер команды}

\subsection{Организация работы -- \PZero}
\begin{tasklist}
  \item Зафиксировать название команды.
  \item Выбрать название продукта.
  \item Проверить название и ник бота: после создания ник изменить нельзя.
  \item Получить у организаторов токен MAX.
  \item Найти точную дату и время дедлайна онлайн-этапа.
  \item Создать Git-репозиторий.
  \item Создать доску задач со статусами \code{Backlog / Ready / In progress / Review / Done}.
  \item Назначить ответственного за каждое направление.
  \item Определить правила code review.
  \item Определить формат веток и коммитов.
  \item Согласовать расписание общих проверок.
  \item Зафиксировать критерии готовности MVP.
  \item Назначить владельца финальной сборки и сдачи.
  \item Создать реестр рисков.
  \item Вести журнал продуктовых и технических решений.
\end{tasklist}

\subsection{Контрольные точки}
\begin{enumerate}
  \item Проблема подтверждена.
  \item Сценарий утверждён.
  \item Технический прототип MAX работает.
  \item Основной сценарий проходит полностью.
  \item Ошибки и пограничные состояния обработаны.
  \item Docker и документация готовы.
  \item Презентация и демонстрация отрепетированы.
  \item Версия зафиксирована и отправлена.
\end{enumerate}
\done{у каждой задачи есть ответственный, срок, приоритет и проверяемый результат.}

\section{Продуктовое исследование}
\owner{product lead; в исследовании участвует вся команда}

\subsection{Целевая аудитория -- \PZero}
\begin{tasklist}
  \item Выбрать одного основного пользователя.
  \item Зафиксировать преподавателя как основной сегмент либо обосновать другой выбор.
  \item Уточнить тип преподавателя: вуз, очные занятия, тип занятия, размер группы и частота занятий.
  \item Описать студента как второго участника процесса.
  \item Определить, кто принимает решение о внедрении: преподаватель, кафедра, деканат или IT-служба.
  \item Определить первый возможный университет, факультет или кафедру для пилота.
\end{tasklist}

\subsection{Исследование проблемы -- \PZero}
\begin{tasklist}
  \item Провести интервью с преподавателями.
  \item Провести короткие интервью со студентами.
  \item Узнать, как посещаемость собирают сейчас.
  \item Получить примеры журналов, таблиц или форм.
  \item Измерить примерное время переклички.
  \item Узнать, как часто возникают ошибки.
  \item Узнать, куда затем переносятся результаты.
  \item Узнать, какие способы обхода или обмана встречаются.
  \item Определить основные неудобства преподавателя.
  \item Определить основные неудобства студента.
  \item Проверить потребность в истории посещений и экспорте.
  \item Найти официальную статистику и надёжные источники.
  \item Отделить подтверждённые факты от гипотез команды.
\end{tasklist}

\subsection{Формулировка проблемы -- \PZero}
\begin{summarybox}
\textbf{Шаблон:} [Пользователь] в ситуации [контекст] хочет [результат], но сталкивается с [барьером], из-за чего [последствие].

\medskip
\textbf{Рабочая формулировка:} преподаватель очного занятия хочет быстро получить достоверный список присутствующих, но вынужден проводить перекличку и вручную переносить результаты, из-за чего теряется учебное время и возникают ошибки.
\end{summarybox}
\begin{tasklist}
  \item Подтвердить каждый элемент формулировки.
  \item Определить главное последствие проблемы.
  \item Зафиксировать, почему существующие способы недостаточны.
  \item Объяснить, почему MAX подходит для решения.
  \item Подготовить список источников.
\end{tasklist}

\subsection{Текущий и будущий процессы -- \PZero}
\subsubsection{As Is}
\begin{tasklist}
  \item Описать создание списка группы, начало переклички и подтверждение присутствия.
  \item Указать место фиксации результата и последующего переноса данных.
  \item Найти ручные действия, ошибки и потери времени.
  \item Указать сервисы, документы и участников процесса.
  \item Описать права доступа к результатам.
\end{tasklist}

\subsubsection{To Be}
\begin{tasklist}
  \item Преподаватель открывает Mini App и выбирает группу.
  \item Запускает сессию и получает временный QR.
  \item Студент сканирует QR и открывает Mini App в MAX.
  \item Backend проверяет пользователя, группу, токен и статус сессии.
  \item Студент получает подтверждение.
  \item Преподаватель видит обновляемый список.
  \item Преподаватель завершает сессию и получает итоговый результат.
\end{tasklist}

\subsection{Ценность, гипотеза и метрики -- \PZero}
\begin{tasklist}
  \item Сформулировать ценность для преподавателя.
  \item Сформулировать ценность для студента.
  \item Сформулировать ценность для университета.
  \item Подготовить продуктовую гипотезу.
  \item Выбрать метрики эффекта.
  \item Зафиксировать исходные значения, если их удалось измерить.
  \item Определить способ проверки метрик на пилоте.
\end{tasklist}
\begin{summarybox}
\textbf{Рабочая гипотеза:} если преподаватель сможет запускать QR-сессию посещаемости внутри MAX, он будет получать готовый список присутствующих быстрее и с меньшим количеством ручных действий.
\end{summarybox}

\section{Формирование MVP}
\owner{product lead и технический лидер}

\subsection{Must Have -- \PZero}
\begin{multicols}{2}
\begin{tasklist}
  \item Авторизация через MAX.
  \item Роли преподавателя и студента.
  \item Тестовая учебная группа.
  \item Привязка студента к группе.
  \item Создание занятия или сессии.
  \item Запуск сбора посещаемости.
  \item Создание временного QR.
  \item Открытие Mini App через QR.
  \item Проверка QR-токена.
  \item Проверка пользователя MAX.
  \item Проверка принадлежности к группе.
  \item Запрет повторной отметки.
  \item Успешная отметка присутствия.
  \item Результат для студента.
  \item Список присутствующих преподавателю.
  \item Завершение сессии.
  \item Обработка основных ошибок.
  \item Работа в MAX Mobile и MAX Web.
\end{tasklist}
\end{multicols}

\subsection{Should Have -- \POne}
\begin{tasklist}
  \item Автоматическая смена QR.
  \item Ручное добавление или удаление отметки преподавателем.
  \item Журнал ручных изменений.
  \item Экспорт CSV.
  \item История занятий.
  \item Повторное открытие активной сессии.
  \item Отображение количества присутствующих.
  \item Состояния «опоздал» и «присутствует».
  \item Защита от слишком частых запросов.
  \item Понятная инструкция для первого запуска.
\end{tasklist}

\subsection{Could Have -- \PTwo}
\begin{tasklist}
  \item Аналитика по занятиям.
  \item Импорт студентов из CSV.
  \item Уведомления преподавателю.
  \item Интеграция с групповым чатом MAX.
  \item Полноценная административная панель.
  \item Дополнительная нативная возможность MAX для платформенного бонуса.
\end{tasklist}

\subsection{Won't Have в хакатонной версии}
\begin{tasklist}
  \item Зафиксировать список функций, которые команда сознательно не делает.
  \item Не добавлять AI только ради презентации.
  \item Не делать распознавание лиц и сложную геолокационную проверку.
  \item Не делать микросервисы.
  \item Не подключать реальную LMS без гарантированного доступа.
  \item Не строить систему управления целым университетом.
\end{tasklist}
\done{один пользовательский сценарий полностью даёт преподавателю итоговый список посещаемости.}

\section{UX/UI}
\owner{frontend/UX-участник}

\subsection{Пользовательские пути -- \PZero}
\begin{tasklist}
  \item Новый и повторно вошедший преподаватель.
  \item Запуск, просмотр и завершение сессии.
  \item Новый студент и успешная отметка.
  \item Повторная отметка.
  \item Просроченный QR.
  \item Студент не найден в группе.
  \item Сессия уже закрыта.
\end{tasklist}

\subsection{Макеты -- \PZero}
\begin{tasklist}
  \item Экран загрузки и главный экран преподавателя.
  \item Выбор группы и создание занятия.
  \item Экран QR и живой список присутствующих.
  \item Подтверждение завершения и итоговый список.
  \item Экран студента до подтверждения и после успеха.
  \item Полный набор экранов ошибок.
\end{tasklist}

\subsection{Состояния и требования -- \PZero}
\begin{tasklist}
  \item Для каждого экрана предусмотреть loading, empty, success, error и disabled.
  \item Обработать отсутствие соединения, повторную попытку, истёкшую сессию и недостаточные права.
  \item Использовать MAX UI и адаптировать интерфейс под мобильный экран.
  \item Проверить веб-версию.
  \item Убрать лишние шаги и сделать основные действия понятными без пояснений.
  \item Сделать крупный читаемый QR и показать время до его обновления.
  \item Давать понятную обратную связь после каждого действия.
  \item Проверить контрастность, размеры элементов и доступность обозначений.
\end{tasklist}

\subsection{UX-проверка -- \POne}
\begin{tasklist}
  \item Собрать кликабельный прототип.
  \item Провести 3--5 UX-тестов.
  \item Зафиксировать проблемы и исправить критичные тупики.
\end{tasklist}

\section{Техническая архитектура}
\owner{технический лидер}

\subsection{Архитектурные задачи -- \PZero}
\begin{tasklist}
  \item Утвердить стек: React, TypeScript, Vite, MAX UI, MAX Bridge, Python, FastAPI, PostgreSQL и Docker Compose.
  \item Выбрать структуру monorepo и разделить frontend, backend и инфраструктуру.
  \item Нарисовать архитектурную схему и описать ответственность компонентов.
  \item Спроектировать модель данных и REST API.
  \item Спроектировать интеграцию с MAX и формат ошибок API.
  \item Определить механизм авторизации.
  \item Определить жизненный цикл сессии посещаемости и QR-токена.
  \item Выбрать обновление списка: polling или WebSocket.
  \item Определить хранение времени и часового пояса.
  \item Определить стратегию миграций и журналирования.
  \item Зафиксировать ключевые решения в коротких ADR.
\end{tasklist}

\subsection{Основные сущности}
\begin{center}
\begin{tabularx}{\textwidth}{*{3}{>{\ttfamily}X}}
users & teachers & students\\
groups & enrollments & lessons\\
attendance\_sessions & qr\_tokens & checkins\\
audit\_log & &\\
\end{tabularx}
\end{center}
\done{любой участник способен объяснить компоненты, связи, движение данных и причины выбора технологий.}

\section{Интеграция с MAX}
\owner{frontend/MAX-разработчик}

\subsection{Обязательная интеграция -- \PZero}
\begin{tasklist}
  \item Получить и безопасно сохранить токен.
  \item Создать бота с окончательным ником и настроить команду \code{/start}.
  \item Подключить Mini App к боту и развернуть его по HTTPS.
  \item Подключить MAX Bridge.
  \item Получать \code{window.WebApp.initData} и \code{start\_param}.
  \item Создавать deep link \code{?startapp=<token>}.
  \item Передавать \code{initData} backend-серверу.
  \item Реализовать серверную HMAC-проверку и проверку свежести \code{auth\_date}.
  \item Проверить запуск по QR, MAX Mobile и MAX Web.
  \item Проверить повторный запуск.
  \item Настроить webhook, если бот получает события.
  \item Предусмотреть режим локальной разработки вне MAX.
  \item Не использовать \code{initDataUnsafe} как доказательство личности.
\end{tasklist}

\subsection{Дополнительное использование платформы -- \POne}
\begin{tasklist}
  \item Добавить полезное сообщение бота после \code{/start}.
  \item Добавить кнопку открытия Mini App и инструкцию преподавателю.
  \item Проверить нативную обратную связь MAX.
  \item Определить возможность для платформенного бонуса.
  \item Описать расширенное использование MAX в README и презентации.
\end{tasklist}

\section{Backend}
\owner{backend-разработчик}

\subsection{Базовая инфраструктура -- \PZero}
\begin{tasklist}
  \item Создать FastAPI-приложение и подключить PostgreSQL.
  \item Настроить SQLAlchemy и Alembic.
  \item Добавить конфигурацию окружения, health-check и единый формат ошибок.
  \item Добавить структурированные логи.
  \item Создать начальные миграции и тестовые данные.
\end{tasklist}

\subsection{Авторизация -- \PZero}
\begin{tasklist}
  \item Создать endpoint авторизации через MAX.
  \item Проверять подпись \code{initData} и срок действия.
  \item Создавать или обновлять пользователя и определять его роль.
  \item Выдавать короткоживущую пользовательскую сессию.
  \item Проверять права для каждого endpoint.
\end{tasklist}

\subsection{Группы -- \PZero}
\begin{tasklist}
  \item Создать тестовую группу.
  \item Получать группы преподавателя и участников группы.
  \item Привязывать студентов к группе.
  \item Проверять доступ преподавателя к группе.
\end{tasklist}

\subsection{Посещаемость -- \PZero}
\begin{tasklist}
  \item Создание и получение активной сессии.
  \item Создание QR-токена, хранение его хеша и срока действия.
  \item Ротация QR и проверка статуса сессии.
  \item Endpoint check-in и атомарная запись отметки.
  \item Запрет дублей через ограничение БД.
  \item Проверка принадлежности студента к группе.
  \item Получение списка отметившихся.
  \item Закрытие сессии и запрет отметок после закрытия.
\end{tasklist}

\subsection{Дополнительные функции -- \POne}
\begin{tasklist}
  \item Ручная корректировка и audit log.
  \item Экспорт CSV и история занятий.
  \item Сводная статистика и импорт CSV.
\end{tasklist}

\subsection{Контракт API -- \PZero}
\begin{tasklist}
  \item Зафиксировать OpenAPI 3.0/3.1.
  \item Проверить status codes и обязательные поля.
  \item Добавить примеры запросов и ответов.
  \item Подготовить тестовые роли и файл \code{DATA-API.yaml}.
\end{tasklist}

\section{Frontend}
\owner{frontend-разработчик}

\subsection{Общая часть -- \PZero}
\begin{tasklist}
  \item Создать React/TypeScript-проект.
  \item Подключить MAX UI, MAX Bridge и API-клиент.
  \item Настроить авторизацию, маршрутизацию по ролям и адаптивную вёрстку.
  \item Добавить обработку загрузки и ошибок.
\end{tasklist}

\subsection{Интерфейс преподавателя -- \PZero}
\begin{tasklist}
  \item Список и выбор групп.
  \item Запуск сессии.
  \item QR с таймером и автоматическим обновлением.
  \item Обновляемый список и счётчик присутствующих.
  \item Завершение сессии и итоговый экран.
\end{tasklist}

\subsection{Интерфейс студента -- \PZero}
\begin{tasklist}
  \item Чтение \code{start\_param} и получение информации о сессии.
  \item Отображение занятия и подтверждение отметки.
  \item Экран успеха.
  \item Состояния повтора, просроченного QR, закрытой сессии и отсутствия в группе.
  \item Восстановление после ошибки соединения.
\end{tasklist}

\subsection{Дополнения -- \POne}
\begin{tasklist}
  \item Ручное исправление списка и экспорт.
  \item История занятий.
  \item Улучшенная анимация и обратная связь.
  \item Первый onboarding преподавателя.
\end{tasklist}

\section{Безопасность и персональные данные}
\owner{backend/DevOps; проверяет вся команда}

\subsection{Обязательные меры -- \PZero}
\begin{tasklist}
  \item Создать модель угроз и определить минимально необходимый набор данных.
  \item Не запрашивать номер телефона без необходимости.
  \item Использовать синтетические данные на демонстрации.
  \item Разграничить права преподавателя и студента.
  \item Валидировать MAX \code{initData} на backend и проверять \code{auth\_date}.
  \item Использовать криптографически случайные QR-токены.
  \item Не помещать персональные данные в QR и хранить только хеш токена.
  \item Ограничить срок действия QR и запретить повторные отметки.
  \item Добавить rate limiting и валидацию входных данных.
  \item Не логировать токены и персональные данные.
  \item Не хранить секреты в коде; добавить \code{.env} в \code{.gitignore} и создать \code{.env.example}.
  \item Проверить историю Git на секреты.
  \item Зафиксировать версии и проверить лицензии зависимостей.
  \item Описать порядок удаления и срок хранения данных.
  \item Подготовить политику обработки данных для прототипа.
  \item Указать соблюдение 152-ФЗ как условие реального внедрения.
\end{tasklist}

\subsection{Антифрод -- \POne}
\begin{tasklist}
  \item Оценить риск пересылки QR.
  \item Ввести ротацию QR и ограничить временное окно.
  \item Показывать преподавателю подозрительные события.
  \item Добавить журнал исправлений.
  \item Не использовать IP или геолокацию как единственное доказательство присутствия.
  \item Честно описать ограничения защиты.
\end{tasklist}

\section{Тестирование}
\owner{QA/DevOps; разработчики пишут автоматические тесты}

\subsection{Unit-тесты -- \PZero}
\begin{tasklist}
  \item Проверка подписи MAX и срока действия.
  \item Генерация и истечение QR-токена.
  \item Проверка ролей, дублей, закрытой сессии и принадлежности к группе.
\end{tasklist}

\subsection{Интеграционные тесты -- \PZero}
\begin{tasklist}
  \item Авторизация MAX $\rightarrow$ пользователь.
  \item Создание сессии $\rightarrow$ QR.
  \item QR $\rightarrow$ check-in $\rightarrow$ обновлённый список.
  \item Завершение $\rightarrow$ запрет новых отметок.
  \item Миграции на пустой базе и повторный запуск приложения.
\end{tasklist}

\subsection{End-to-end -- \PZero}
\begin{tasklist}
  \item Полный путь преподавателя и полный путь студента.
  \item Два разных студента и повторная отметка.
  \item Просроченный QR, неизвестный студент и закрытая сессия.
  \item Потеря соединения, ошибка backend и восстановление после ошибки.
\end{tasklist}

\subsection{Совместимость -- \PZero}
\begin{tasklist}
  \item MAX Android, MAX iOS при наличии и MAX Web.
  \item Разные размеры экрана и медленное соединение.
  \item Повторное открытие Mini App.
\end{tasklist}

\subsection{Нагрузочная проверка -- \POne}
\begin{tasklist}
  \item Одновременная отметка всей группы.
  \item 30--100 запросов check-in за короткий период.
  \item Отсутствие дублей при параллельных запросах.
  \item Измерение времени ответа API и поведения PostgreSQL.
\end{tasklist}
\done{основной сценарий стабильно проходит несколько раз подряд без ручного восстановления среды.}

\section{Docker, DevOps и развёртывание}
\owner{DevOps/QA}

\subsection{Воспроизводимый запуск -- \PZero}
\begin{tasklist}
  \item Dockerfile для frontend и backend.
  \item PostgreSQL и все локальные компоненты в \code{compose.yaml}.
  \item \code{.dockerignore}, \code{.env.example} и зафиксированные версии образов.
  \item Health-check компонентов, миграции и тестовые данные при запуске.
  \item Одна команда запуска и одна команда остановки.
  \item Проверка запуска на чистой машине и повторного запуска.
  \item Проверка сборки менее чем за 5 минут без загрузки базовых образов.
  \item HTTPS, домен, постоянно доступные API и Mini App.
  \item Логи, минимальный мониторинг и резервная копия БД.
  \item Проверка отсутствия секретов в образах.
\end{tasklist}

\subsection{Автоматизация -- \POne}
\begin{tasklist}
  \item CI для тестов и Docker build.
  \item Автоматический deploy.
  \item Проверка зависимостей и secret scanning.
\end{tasklist}

\section{Техническая документация}
\owner{технический лидер; проверяет участник без контекста}

\subsection{README -- \PZero}
\begin{tasklist}
  \item Назначение проекта, решаемая проблема и основной сценарий.
  \item Архитектура и состав компонентов.
  \item Требования к окружению, переменные и порты.
  \item Одна команда запуска, команда остановки и повторный запуск.
  \item Настройка MAX и работа с тестовыми данными.
  \item Тестовые роли и пошаговый сценарий проверки.
  \item Ожидаемый результат.
  \item Внешние сервисы, описание данных и известные ограничения.
  \item Безопасность, секреты и лицензии.
\end{tasklist}

\subsection{API-документация -- \PZero}
\begin{tasklist}
  \item \code{openapi.yaml} или \code{openapi.json} версии 3.0/3.1.
  \item Примеры запросов, ответов, ошибок и status codes.
  \item Тестовые аккаунты и тестовый JSON/CSV.
  \item \code{DATA-API.yaml} и перечень обязательных проверок API.
\end{tasklist}

\section{Пилот и масштабирование}
\owner{product lead}

\subsection{Пилот -- \POne}
\begin{tasklist}
  \item Выбрать первый вуз, факультет или кафедру.
  \item Определить количество преподавателей, групп и продолжительность пилота.
  \item Описать подключение пользователей и обучение преподавателя.
  \item Определить необходимые данные, владельца процесса и канал поддержки.
  \item Выбрать метрики и критерии успеха.
  \item Описать действия после успешного и отрицательного результата.
\end{tasklist}

\subsection{Масштабирование -- \POne}
\begin{tasklist}
  \item Определить ядро продукта и переменную часть.
  \item Выбрать следующий факультет и следующий университет.
  \item Описать интеграции, адаптацию модели данных и организационные ресурсы.
  \item Описать юридические ограничения и риски тиражирования.
  \item Составить реалистичную последовательность масштабирования.
\end{tasklist}

\section{Презентация}
\owner{product lead; участвует вся команда}

\subsection{Служебный слайд -- \PZero}
\begin{tasklist}
  \item Ссылка на работающий продукт в MAX.
  \item Ссылка на Git-репозиторий и commit hash.
  \item HTTPS-адрес API.
  \item Тестовые аккаунты и проверочные переменные окружения.
  \item Краткий порядок проверки.
\end{tasklist}

\subsection{Продуктовая часть -- \PZero}
\begin{tasklist}
  \item Название, команда и executive summary.
  \item Целевая аудитория и подтверждённая проблема.
  \item Источники и результаты интервью.
  \item As Is, To Be, решение и основной сценарий.
  \item Ценность, метрики и преимущество перед текущим процессом.
\end{tasklist}

\subsection{Техническая часть -- \PZero}
\begin{tasklist}
  \item Архитектура, интеграция с MAX и модель данных.
  \item Технологии, безопасность и типы используемых данных.
  \item Ограничения MVP и известные риски.
\end{tasklist}

\subsection{Развитие -- \POne}
\begin{tasklist}
  \item План пилота, масштабирование и необходимые интеграции.
  \item Следующие шаги и условия внедрения.
\end{tasklist}

\section{Демонстрация и защита}
\owner{лидер презентации; участвует вся команда}

\subsection{Подготовка -- \PZero}
\begin{tasklist}
  \item Написать сценарий демонстрации.
  \item Подготовить роли преподавателя и студента, тестовую группу и устройства.
  \item Подготовить запасную сеть и резервное видео демонстрации.
  \item Проверить ссылки, MAX и API перед выступлением.
  \item Подготовить короткое выступление и распределить слайды.
  \item Подготовить ответы на вопросы.
  \item Отрепетировать защиту с таймером и провести стресс-репетицию.
  \item Убедиться, что каждый участник понимает весь продукт.
  \item Разделять в ответах факты, гипотезы и ограничения.
\end{tasklist}

\subsection{Вероятные вопросы жюри}
\begin{itemize}
  \item Почему нельзя использовать обычную форму или таблицу?
  \item Что мешает студенту отправить QR другу?
  \item Почему выбран MAX и кто основной пользователь?
  \item Чем подтверждена проблема и как измеряется результат?
  \item Какие персональные данные хранит система?
  \item Как решение интегрируется с университетом и работает при плохом интернете?
  \item Почему выбран этот стек и как система масштабируется?
  \item Где провести первый пилот и что не вошло в MVP?
\end{itemize}

\section{Финальная сдача}
\owner{лидер команды и DevOps}

\subsection{Финальный чек-лист -- \PZero}
\begin{tasklist}
  \item Провести финальный code review и удалить отладочный код.
  \item Проверить Git на секреты и зависимости на допустимые лицензии.
  \item Запустить все тесты и собрать Docker с нуля.
  \item Пройти основной сценарий строго по README.
  \item Проверить OpenAPI, DATA-API, тестовые аккаунты и тестовые данные.
  \item Проверить публичный HTTPS, бота, Mini App, мобильную и веб-версию.
  \item Создать финальный commit и зафиксировать commit hash.
  \item Поставить Git tag.
  \item Создать архив при необходимости и рассчитать контрольную сумму.
  \item Экспортировать презентацию в PDF.
  \item Проверить служебный слайд и все ссылки.
  \item Отправить материалы и больше не изменять зафиксированную версию.
\end{tasklist}

\section{Распределение задач на четырёх участников}

\begin{longtable}{>{\raggedright\arraybackslash}p{11mm}>{\raggedright\arraybackslash}p{42mm}>{\raggedright\arraybackslash}p{52mm}>{\raggedright\arraybackslash}p{52mm}}
\toprule
\textbf{№} & \textbf{Роль} & \textbf{Основная зона} & \textbf{Дополнительная зона}\\
\midrule
\endfirsthead
\toprule
\textbf{№} & \textbf{Роль} & \textbf{Основная зона} & \textbf{Дополнительная зона}\\
\midrule
\endhead
1 & Product Lead & Исследование, ценность, пилот, масштабирование & Презентация и защита\\
2 & Frontend/MAX & React, MAX UI, MAX Bridge, UX & Интеграционные тесты\\
3 & Backend & FastAPI, PostgreSQL, check-in, безопасность & OpenAPI и DATA-API\\
4 & QA/DevOps & Docker, deploy, тестирование, документация & Демонстрация и финальная сборка\\
\bottomrule
\end{longtable}

\begin{summarybox}
Code review, проверка основного сценария, подготовка к вопросам и понимание продукта являются общей ответственностью команды. Разделение ролей не должно создавать изолированные зоны знаний.
\end{summarybox}

\section{Последовательность выполнения}

\begin{center}
\begin{tcolorbox}[width=0.82\textwidth,colback=Panel,colframe=MaxCyan,arc=3mm,boxrule=1pt]
\centering
\textbf{Исследование проблемы}\\[1mm]
$\downarrow$\\
\textbf{Целевая аудитория и метрики}\\[1mm]
$\downarrow$\\
\textbf{As Is / To Be}\\[1mm]
$\downarrow$\\
\textbf{Границы MVP}\\[1mm]
$\downarrow$\\
\textbf{UX-прототип}\\[1mm]
$\downarrow$\\
\textbf{Технический spike MAX}\\[1mm]
$\downarrow$\\
\textbf{Backend + Frontend}\\[1mm]
$\downarrow$\\
\textbf{Интеграция основного сценария}\\[1mm]
$\downarrow$\\
\textbf{Безопасность и обработка ошибок}\\[1mm]
$\downarrow$\\
\textbf{Тестирование}\\[1mm]
$\downarrow$\\
\textbf{Docker и документация}\\[1mm]
$\downarrow$\\
\textbf{Презентация и демонстрация}\\[1mm]
$\downarrow$\\
\textbf{Фиксация версии и сдача}
\end{tcolorbox}
\end{center}

\clearpage
\section{Шаблон еженедельного статуса}

\begin{longtable}{>{\bfseries\raggedright\arraybackslash}p{48mm}>{\raggedright\arraybackslash}p{112mm}}
\toprule
Период & \rule{0pt}{9mm}\\
Главная цель недели & \rule{0pt}{9mm}\\
Готово & \rule{0pt}{9mm}\\
В работе & \rule{0pt}{9mm}\\
Заблокировано & \rule{0pt}{9mm}\\
Риски & \rule{0pt}{9mm}\\
Решения команды & \rule{0pt}{9mm}\\
Метрика недели & \rule{0pt}{9mm}\\
Демонстрируемый результат & \rule{0pt}{9mm}\\
Следующая контрольная точка & \rule{0pt}{9mm}\\
\bottomrule
\end{longtable}

\vfill
\begin{center}
\color{Muted}\small
Документ является рабочим чек-листом. Перед сдачей требования следует повторно сверить с актуальными правилами хакатона и документацией MAX.
\end{center}

\end{document}
